\documentclass[11pt]{amsart}

\usepackage[a4paper,margin=2.8cm]{geometry}
\usepackage[T1]{fontenc}
\usepackage{lmodern}
\usepackage{microtype}
\usepackage{amsmath,amssymb,mathtools}
\usepackage{enumitem}
\usepackage{xcolor}
\usepackage{hyperref}

\hypersetup{
  colorlinks=true,
  linkcolor=blue!45!black,
  citecolor=blue!45!black,
  urlcolor=blue!45!black
}
\setlist[itemize]{leftmargin=*,itemsep=0.25em,topsep=0.4em}
\setlist[enumerate]{leftmargin=*,itemsep=0.25em,topsep=0.4em}

\newtheorem{theorem}{Theorem}[section]
\newtheorem{proposition}[theorem]{Proposition}
\newtheorem{lemma}[theorem]{Lemma}
\newtheorem{corollary}[theorem]{Corollary}
\theoremstyle{definition}
\newtheorem{definition}[theorem]{Definition}
\theoremstyle{remark}
\newtheorem{remark}[theorem]{Remark}

\newcommand{\M}{M}
\newcommand{\Met}{\operatorname{Met}}
\newcommand{\dMet}{d_{\ell_M}}
\newcommand{\TV}{\operatorname{TV}}
\newcommand{\Per}{\operatorname{Per}}
\newcommand{\Vol}{\operatorname{Vol}}
\newcommand{\BV}{\operatorname{BV}}
\newcommand{\Lip}{\operatorname{Lip}}
\newcommand{\esssup}{\operatorname*{ess\,sup}}
\newcommand{\dmu}[1]{\,d\mu_{#1}}
\newcommand{\one}{\mathbf 1}
\newcommand{\R}{\mathbb R}
\newcommand{\cI}{\mathcal I}
\newcommand{\cS}{\mathcal S}
\newcommand{\cE}{\mathcal E}
\newcommand{\hch}{\mathsf h}

\title[Rough isoperimetry]{Rough isoperimetry\\
\large A direct programme and first theorems}
\author{Glen Wheeler}
\author{Lashi Bandara}
\date{11 August 2026}

\begin{document}

\begin{abstract}
We propose a short route from the structure theory of rough Riemannian
metrics to geometric variational problems.  There is one important codimension-one subtlety: a
measurable metric is defined only up to sets of ambient measure zero, and hence
cannot naively be restricted to a hypersurface.  We use instead the
$L^1$-relaxation of the metric total variation, in the spirit of
Amar--Bellettini.  On a compact manifold we prove existence of isoperimetric
regions of every prescribed volume, continuity of the rough isoperimetric
profile (with local H\"older regularity), and compactness and convergence of
minimisers under convergence in the Bandara--Hassan metric on the space of
rough metrics.  The normalised profiles obey a locally uniform
$(n-1)/n$-H\"older stability estimate.  We also prove a
rough Cheeger inequality and explicit log-Lipschitz estimates for the
isoperimetric constant, the Cheeger constant, and every positive eigenvalue of
the rough Laplacian.  This is, I feel, a substantial package of theorems.
\end{abstract}

\maketitle

\section{The recommendation}

Consider the isoperimetric problem
\begin{equation}\label{eq:direct-problem}
 \min\bigl\{\Per_g(E):\mu_g(E)=v\bigr\}.
\end{equation}
For smooth $g$, the regular boundary of a minimiser has constant mean
curvature.  
 For a
merely measurable $g$, the equation $H=\lambda$ is not meaningful pointwise, but
\eqref{eq:direct-problem} remains meaningful and gives a natural
definition of a \emph{variational rough CMC analogue}.

Lashi's observations make this particularly nice.
A rough metric gives a
Radon measure, common $L^p$ and Sobolev spaces throughout each component of
the space of rough metrics, a closed Dirichlet form and Laplacian, and a
canonical intrinsic distance.  More recently, the collection of rough
metrics itself has been made into a complete length space
\cite{BandaraHassan2025}.  On a compact manifold every pair of rough metrics
lies at finite distance.  The distance is exactly strong enough to compare
volume, codimension-one energy, and Dirichlet energy with explicit constants.

The proposed short paper therefore has three layers.

\begin{itemize}
\item Define a canonical relaxed perimeter and solve
  \eqref{eq:direct-problem} for every $v$.
\item Prove that the entire isoperimetric problem is stable as the rough
  metric moves in the Bandara--Hassan space, quantitatively on interior volume
  fractions.
\item Join this geometric side to Lashi's operator-theoretic side through a
  rough Cheeger inequality and quantitative spectral stability.
\end{itemize}

I have had a go at all three layers below.  The main remaining problem is regularity: how much regularity of $g$ is needed to
upgrade this variational analogue to a recognisable hypersurface satisfying a
weak or classical CMC equation?

The direct method and the abstract Cheeger mechanism are classical; we do not
claim otherwise.  The short-paper contribution proposed here is their unified
formulation for almost-everywhere measurable metrics in the
Bandara--Hassan topology, the sharp codimension-one comparison, and the joint
stability of profiles and minimisers with quantitative constants.
For a smooth-metric profile-stability argument based on perimeter comparison,
compactness, and volume repair, see also Zheng \cite{Zheng2020}.  The point here
is the extension to relaxed perimeter for measurable coefficients and the
$d_{\ell_M}$ topology.

\section{A codimension-one point to fix}

A rough metric is an almost-everywhere equivalence class.  A hypersurface has
zero $n$-dimensional measure.  Consequently, the expression ``the metric
induced by $g$ on the hypersurface'' is not determined by the rough metric:
one could change $g$ on that hypersurface without changing the rough metric at
all.

There is a related functional-analytic issue.
If rough mass is defined by testing a current only against \emph{smooth}
forms satisfying an almost-everywhere comass bound, the smooth test class can
fail to see the discontinuities of the coefficients.  Amar and Bellettini
identified exactly this phenomenon for discontinuous Finsler metrics; they
showed in particular that smooth test fields may be inadequate
\cite{AmarBellettini1994}.  Their remedy is a lower-semicontinuous generalised
total variation, equivalently described by relaxation and by a larger class
of divergence-controlled fields.
This enables more straightforward existence proofs.

\section{Rough metrics and relaxed perimeter}

Let $M^n$ be a smooth, compact, connected manifold without boundary, with
$n\geq2$.  Orientation is not needed.  A \emph{rough Riemannian metric} $g$ is
a measurable positive-definite symmetric two-tensor such that, relative to
one (and hence every) smooth background metric, both $g$ and $g^{-1}$ are
essentially bounded.  Equivalently, $g$ is locally bounded and locally
elliptic.  It induces the Radon measure $\mu_g$ and a measurable dual norm
$|\xi|_{g^*}$ on covectors.

For $g,\widetilde g\in\Met(M)$, write
\begin{equation}\label{eq:metric-on-metrics}
 \dMet(g,\widetilde g)
 :=\inf\left\{\log C:
 C^{-1}|X|_g\leq |X|_{\widetilde g}\leq C|X|_g
 \text{ a.e. for every }X\in TM\right\}.
\end{equation}
This is the Bandara--Hassan metric, with notation slightly shortened.  Since
$M$ is compact, it is finite for every pair of rough metrics
\cite{BandaraHassan2025}.

For $u\in L^1(M,\mu_g)$ define
\begin{equation}\label{eq:rough-TV}
 \TV_g(u):=\inf\left\{
 \liminf_{j\to\infty}\int_M |du_j|_{g^*}\dmu g:
 u_j\in C^\infty(M),\quad u_j\to u\text{ in }L^1(M,\mu_g)
 \right\}.
\end{equation}
For a measurable set $E\subset M$, its relaxed rough perimeter is
\begin{equation}\label{eq:rough-perimeter}
 \Per_g(E):=\TV_g(\one_E).
\end{equation}
This depends only on the almost-everywhere class of $g$, is lower
semicontinuous under $L^1$ convergence, and agrees with ordinary Riemannian
perimeter when $g$ is smooth.  In local charts it is the relaxed total
variation associated with a bounded measurable elliptic Finsler integrand;
this is the setting developed in
\cite[Theorem~5.1]{AmarBellettini1994} and
\cite{BellettiniBouchitteFragala1999}.  Indeed, fix a smooth background metric
$h$, put $\rho_{g,h}=d\mu_g/d\mu_h$, and set
\[
 F_{g,h}(x,\xi):=\rho_{g,h}(x)|\xi|_{g^*}.
\]
There are constants $0<c\leq C<\infty$ such that
\[
 c|\xi|_{h^*}\leq F_{g,h}(x,\xi)\leq C|\xi|_{h^*},
 \qquad
 |F_{g,h}(x,\xi)-F_{g,h}(x,\zeta)|
 \leq C|\xi-\zeta|_{h^*}
\]
almost everywhere.  Smooth density in $W^{1,1}(M,h)$ therefore shows that the
infimum is unchanged if the approximants are taken in $W^{1,1}$, and permits
chartwise localisation and partition of unity.  A precise dual formula uses
the enlarged test-field class of Amar--Bellettini, rather than only smooth
fields; none of the results below needs that formula.  From now on, unlabelled
$L^1$ convergence means convergence with respect to $\mu_h$ for one fixed
smooth $h$; uniform ellipticity makes this equivalent to convergence for any
rough metric under discussion.

The next elementary linear-algebra estimate is key.

\begin{lemma}[The cofactor estimate]\label{lem:cofactor}
Let $g,\widetilde g\in\Met(M)$ and set
$\delta=\dMet(g,\widetilde g)$.  For every $p\geq1$ and every measurable
covector field $\xi$,
\begin{align}
 e^{-n\delta}\,d\mu_g
 &\leq d\mu_{\widetilde g}\leq e^{n\delta}\,d\mu_g,
 \label{eq:measure-comparison}\qquad\\
 e^{-(n+p-2)\delta}|\xi|_{g^*}^p\,d\mu_g
 &\leq |\xi|_{\widetilde g^*}^p\,d\mu_{\widetilde g}
 \leq e^{(n+p-2)\delta}|\xi|_{g^*}^p\,d\mu_g.
 \label{eq:p-energy-comparison}
\end{align}
The inequalities are understood as inequalities of measurable densities.
\end{lemma}

\begin{proof}
Fix $d>\delta$; at the end let $d\downarrow\delta$.  At almost every point,
choose a
$g$-orthonormal basis diagonalising $\widetilde g$ and write
\[
 |X|_{\widetilde g}^2=\sum_{i=1}^n s_i^2(X^i)^2,
 \qquad e^{-d}\leq s_i\leq e^d.
\]
Then
\[
 \frac{d\mu_{\widetilde g}}{d\mu_g}=\prod_{i=1}^n s_i,
 \qquad
 |\xi|_{\widetilde g^*}^2=\sum_{i=1}^n s_i^{-2}\xi_i^2.
\]
The first assertion follows immediately.  If
$s_-:=\min_i s_i$ and $s_+:=\max_i s_i$, then
\[
 \left(\prod_i s_i\right)s_+^{-p}|\xi|_{g^*}^p
 \leq
 \left(\prod_i s_i\right)|\xi|_{\widetilde g^*}^p
 \leq
 \left(\prod_i s_i\right)s_-^{-p}|\xi|_{g^*}^p.
\]
One factor in the product is $s_+$ and the other $n-1$ factors are at
least $e^{-d}$; similarly, one factor is $s_-$ and the other factors
are at most $e^d$.  Since $p\geq1$,
\[
 \left(\prod_i s_i\right)s_+^{-p}\geq e^{-(n+p-2)d},
 \qquad
 \left(\prod_i s_i\right)s_-^{-p}\leq e^{(n+p-2)d}.
\]
Letting $d\downarrow\delta$ proves the result.
\end{proof}

The exponent $n-1$ at $p=1$ is precisely the codimension-one exponent; the
exponent $n$ at $p=2$ controls the Dirichlet form.  The general exponent is
pointwise sharp: in a $g$-orthonormal frame take
$s_1=e^{-\delta}$, $s_2=\cdots=s_n=e^\delta$, and let $\xi$ point in the
first coordinate direction.  The upper density ratio is then exactly
$e^{(n+p-2)\delta}$.

\begin{proposition}[Uniform comparison of rough perimeters]
\label{prop:perimeter-comparison}
Let $g,\widetilde g\in\Met(M)$ and set
$\delta=\dMet(g,\widetilde g)$.  For every measurable $E\subset M$,
\begin{equation}\label{eq:perimeter-comparison}
 e^{-(n-1)\delta}\Per_g(E)
 \leq \Per_{\widetilde g}(E)
 \leq e^{(n-1)\delta}\Per_g(E).
\end{equation}
In particular, the class of finite-perimeter sets is independent of the rough
metric on compact $M$.  It is the usual class of Caccioppoli sets for any
smooth background metric.  Every sequence of characteristic functions with
uniformly bounded $g$-perimeter has an $L^1$-convergent subsequence;
equivalently, perimeter sublevels are relatively compact.
\end{proposition}

\begin{proof}
The $L^1(M,\mu_g)$ and $L^1(M,\mu_{\widetilde g})$ topologies coincide by
\eqref{eq:measure-comparison}.  Apply
\eqref{eq:p-energy-comparison} with $p=1$ to every smooth approximating
sequence in \eqref{eq:rough-TV}, and then reverse the roles of the metrics.
Comparison with a fixed smooth metric identifies the finite-perimeter class
and transfers the standard BV compactness theorem \cite{AFP2000}.
\end{proof}

We shall also use the lattice inequality
\begin{equation}\label{eq:lattice}
 \Per_g(E\cup F)+\Per_g(E\cap F)
 \leq \Per_g(E)+\Per_g(F).
\end{equation}
For $u,v\in W^{1,1}$, the gradients of $u\vee v$ and $u\wedge v$ agree almost
everywhere with $du,dv$ in the two possible orders.  Hence
\[
 F_{g,h}(x,d(u\vee v))+F_{g,h}(x,d(u\wedge v))
 =F_{g,h}(x,du)+F_{g,h}(x,dv)
\]
almost everywhere.  Choose joint smooth near-recovery sequences for
$\one_E$ and $\one_F$, apply this identity, approximate the resulting maximum
and minimum strongly in $W^{1,1}(M,h)$, and take a diagonal sequence.  The
Lipschitz estimate for $F_{g,h}$ above and lower semicontinuity give
\eqref{eq:lattice}.  In particular,
\begin{equation}\label{eq:union-cost}
 \Per_g(E\cup F)\leq\Per_g(E)+\Per_g(F),
 \qquad
 \Per_g(E\setminus F)\leq\Per_g(E)+\Per_g(F).
\end{equation}

\section{Exact-volume rough isoperimetry}

Write $V_g:=\mu_g(M)$.  Define the rough isoperimetric profile by
\begin{equation}\label{eq:profile}
 \cI_g(v):=\inf\bigl\{\Per_g(E):E\subset M\text{ measurable},\
 \mu_g(E)=v\bigr\},
 \qquad 0\leq v\leq V_g.
\end{equation}

\begin{lemma}[The constraint class is nonempty]\label{lem:constraint}
For every $v\in[0,V_g]$ there is a finite-perimeter set $E$ with
$\mu_g(E)=v$.
\end{lemma}

\begin{proof}
Take a finite smooth triangulation, subdivided to be subordinate to a finite
coordinate atlas, and order its $n$-simplices.  Their
boundaries have zero $\mu_g$-measure.  Add simplex interiors until the target
volume lies in the next simplex.  In a smooth coordinate chart on that
simplex, sweep the remaining part by parallel coordinate hyperplanes.  The
enclosed $\mu_g$-volume is a continuous function of the sweep parameter: the
density of $\mu_g$ is bounded and positive relative to smooth volume, and
each slicing hyperplane is null.  The intermediate value theorem gives the
exact remaining volume.  The resulting set has piecewise smooth boundary and
hence finite perimeter for a smooth background metric, and therefore for $g$
by Proposition~\ref{prop:perimeter-comparison}.
\end{proof}

\begin{theorem}[Existence of rough isoperimetric regions]
\label{thm:existence}
For every $v\in[0,V_g]$ there is a finite-perimeter set $E_v$ such that
\[
 \mu_g(E_v)=v,
 \qquad
 \Per_g(E_v)=\cI_g(v).
\]
Moreover, $\cI_g(v)=\cI_g(V_g-v)$, and $\cI_g(v)>0$ for
$0<v<V_g$.
\end{theorem}

\begin{proof}
The endpoint cases are immediate, and Lemma~\ref{lem:constraint} makes the
infimum finite.  Let $E_j$ be a minimising sequence at an interior volume.
Its $g$-perimeters are bounded.  Proposition~\ref{prop:perimeter-comparison}
and BV compactness for a fixed smooth background metric give, after passage
to a subsequence,
\[
 \one_{E_j}\longrightarrow\one_E\quad\text{in }L^1(M).
\]
The density of $\mu_g$ relative to smooth volume is bounded, so
$\mu_g(E)=v$.  Lower semicontinuity of \eqref{eq:rough-TV} gives
\[
 \Per_g(E)\leq\liminf_{j\to\infty}\Per_g(E_j)=\cI_g(v).
\]
Thus $E$ is a minimiser.  Complementation preserves perimeter and interchanges
$v$ with $V_g-v$, proving symmetry.

For positivity, fix a smooth metric $h$ and let
$D=\dMet(g,h)$.  The smooth relative isoperimetric inequality and
\eqref{eq:measure-comparison}--\eqref{eq:perimeter-comparison} give
\begin{equation}\label{eq:rough-relative-iso}
 \Per_g(E)\geq c_h e^{-2(n-1)D}
 \min\{\mu_g(E),\mu_g(M\setminus E)\}^{\frac{n-1}{n}}.
\end{equation}
This is positive at a nontrivial prescribed volume.
\end{proof}

The next lemma is the small device that yields both continuity of the profile
and stability under change of metric.

\begin{lemma}[Uniform repair of a small volume error]
\label{lem:uniform-repair}
Fix a smooth metric $h_0$ on $M$, and let $D,\alpha>0$.  Set
\[
 \mathcal B_D(h_0):=\{q\in\Met(M):\dMet(q,h_0)\leq D\},
 \qquad \gamma:=\frac{n-1}{n}.
\]
There are constants $s_0=s_0(M,h_0,\alpha)>0$ and
$C=C(M,h_0,D,\alpha)<\infty$ with the following property.  If
$q\in\mathcal B_D(h_0)$ and $E$ has finite $q$-perimeter with
\[
 \min\{\mu_q(E),\mu_q(M\setminus E)\}\geq\alpha,
\]
then, for every $s\in\R$ with $|s|\leq s_0$, there is a finite-perimeter set
$E_s$ such that
\begin{equation}\label{eq:uniform-repair}
 \mu_q(E_s)=\mu_q(E)+s,
 \qquad
 \Per_q(E_s)\leq\Per_q(E)+C|s|^\gamma,
 \qquad
 \mu_q(E_s\mathbin\triangle E)=|s|.
\end{equation}
The constants are uniform in both $q$ and $E$.
\end{lemma}

\begin{proof}
Choose $r_0>0$ smaller than the injectivity radius of $(M,h_0)$.  By
compactness, after decreasing $r_0$ if necessary, there are $a_0,b_0>0$ such
that, for every $x\in M$ and $0<r<r_0$,
\[
 \mu_{h_0}(B_{h_0}(x,r))\geq a_0r^n,
 \qquad
 \Per_{h_0}(B_{h_0}(x,r))\leq b_0r^{n-1}.
\]
Put $\kappa:=a_0/\mu_{h_0}(M)$.  If $A\subset M$ is measurable and
$\mu_q(A)\geq\alpha$, Fubini's theorem gives
\begin{align*}
 \int_M\mu_q(A\cap B_{h_0}(x,r))\dmu{h_0}
 &=\int_A\mu_{h_0}(B_{h_0}(y,r))\dmu q\\
 &\geq a_0r^n\mu_q(A).
\end{align*}
Consequently, for some $x=x(A,r)$,
\begin{equation}\label{eq:averaging-ball}
 \mu_q(A\cap B_{h_0}(x,r))\geq\kappa\alpha r^n.
\end{equation}
Take
\[
 s_0:=\min\left\{\frac{\alpha}{2},
       \kappa\alpha\left(\frac{r_0}{2}\right)^n\right\}.
\]
Suppose first that $0<s\leq s_0$, set
$r=(s/(\kappa\alpha))^{1/n}\leq r_0/2$, and apply
\eqref{eq:averaging-ball} with $A=M\setminus E$.  For the resulting centre
$x$, the function
\[
 \rho\longmapsto
 \mu_q((M\setminus E)\cap B_{h_0}(x,\rho))
\]
is continuous on $[0,r]$: every $h_0$-geodesic sphere of radius below $r_0$
is smooth and hence $\mu_q$-null.  It is zero at $\rho=0$ and at least $s$ at
$\rho=r$.  Thus some $\rho\leq r$ makes it equal to $s$.  Set
$E_s=E\cup B_{h_0}(x,\rho)$.  The lattice inequality and the perimeter
comparison give
\begin{align*}
 \Per_q(E_s)
 &\leq\Per_q(E)+\Per_q(B_{h_0}(x,\rho))\\
 &\leq\Per_q(E)+e^{(n-1)D}b_0\rho^{n-1}\\
 &\leq\Per_q(E)+e^{(n-1)D}b_0
       (\kappa\alpha)^{-\gamma}s^\gamma.
\end{align*}
For $-s_0\leq s<0$, apply the same argument with $A=E$ and set
$E_s=E\setminus B_{h_0}(x,\rho)$.  The symmetric-difference identity is
immediate from the construction, and the case $s=0$ is trivial.
\end{proof}

\begin{proposition}[Continuity and local H\"older regularity]
\label{prop:profile-continuity}
The function $\cI_g:[0,V_g]\to[0,\infty)$ is continuous.  More precisely,
for every compact interval $K\Subset(0,V_g)$ there is $C_K<\infty$ such that
\begin{equation}\label{eq:profile-holder}
 |\cI_g(v)-\cI_g(w)|\leq C_K|v-w|^{(n-1)/n}
 \qquad(v,w\in K).
\end{equation}
\end{proposition}

\begin{proof}
Fix a smooth metric $h$ on $M$.
For fixed $x$ and all small radii, the function
$r\mapsto\mu_g(B_h(x,r))$ is continuous and strictly increasing, while
$\Per_g(B_h(x,r))\leq Cr^{n-1}$.  Exact-volume balls therefore show that
$\cI_g(v)\to0$ as $v\downarrow0$; symmetry handles the other endpoint.  Now
write $K=[a,b]$ and set $\alpha=\min\{a,V_g-b\}>0$.  Apply
Lemma~\ref{lem:uniform-repair} with $h_0=h$, $D=\dMet(g,h)$, and $q=g$.
If $|v-w|\leq s_0$, repair a
$g$-isoperimetric region of volume $v$ by $s=w-v$ to obtain
\[
 \cI_g(w)\leq\cI_g(v)+C|v-w|^{(n-1)/n}.
\]
Interchanging $v$ and $w$ gives the reverse inequality.  For arbitrary
$v,w\in K$, divide the interval between them into $m$ equal pieces, where
$m\leq1+V_g/s_0$ and each piece has length at most $s_0$.  Summing the local
estimate gives
\[
 |\cI_g(v)-\cI_g(w)|
 \leq C m^{1/n}|v-w|^{(n-1)/n},
\]
which proves \eqref{eq:profile-holder} and interior continuity.
\end{proof}

\section{Stability in the space of rough metrics}

The following is the main geometric payoff of using Lashi's metric on the
space of metrics.

\begin{theorem}[Stability of rough isoperimetry]
\label{thm:profile-stability}
Suppose $g_j,g\in\Met(M)$ and
\[
 \dMet(g_j,g)\longrightarrow0.
\]
Let $v_j\to v$ with $0<v<V_g$ and $0<v_j<V_{g_j}$.  Then
\begin{equation}\label{eq:profile-convergence}
 \cI_{g_j}(v_j)\longrightarrow\cI_g(v).
\end{equation}
If $E_j$ is a $g_j$-isoperimetric region of volume $v_j$, then a subsequence
converges in $L^1$ to a $g$-isoperimetric region $E$ of volume $v$, and
\begin{equation}\label{eq:minimiser-perimeter-convergence}
 \Per_{g_j}(E_j)\longrightarrow\Per_g(E).
\end{equation}
In particular, for each fixed volume fraction $\theta\in(0,1)$,
\[
 \cI_{g_j}(\theta V_{g_j})\longrightarrow\cI_g(\theta V_g).
\]
\end{theorem}

\begin{proof}
Let $E$ be a $g$-isoperimetric region of volume $v$.  By
\eqref{eq:measure-comparison},
\[
 s_j:=v_j-\mu_{g_j}(E)\longrightarrow0.
\]
Fix a smooth $h_0$.  For all large $j$, the metrics $g_j$ lie in one ball
$\mathcal B_D(h_0)$ and both phases of $E$ have $g_j$-volume at least some
$\alpha>0$.  Lemma~\ref{lem:uniform-repair}, applied with $q=g_j$, gives
$F_j$ of exact $g_j$-volume $v_j$ and
\[
 \Per_{g_j}(F_j)\leq\Per_{g_j}(E)+o(1).
\]
Proposition~\ref{prop:perimeter-comparison} implies
$\Per_{g_j}(E)\to\Per_g(E)$.  Consequently,
\begin{equation}\label{eq:stability-limsup}
 \limsup_{j\to\infty}\cI_{g_j}(v_j)\leq\cI_g(v).
\end{equation}

For the matching lower bound, pass first to a subsequence realising
$\liminf_j\cI_{g_j}(v_j)$, and take $g_j$-isoperimetric regions $E_j$.  The
preceding upper bound and \eqref{eq:perimeter-comparison} bound
$\Per_g(E_j)$.  BV compactness gives, after a further subsequence,
$\one_{E_j}\to\one_E$ in $L^1$.  The estimate
\[
 \bigl|\mu_{g_j}(A)-\mu_g(A)\bigr|
 \leq (e^{n\delta_j}-1)\mu_g(A),
 \qquad \delta_j:=\dMet(g_j,g),
\]
is uniform over measurable sets $A$.  Together with $L^1$ convergence it
therefore gives
\[
 \mu_g(E)=\lim_{j\to\infty}\mu_g(E_j)
 =\lim_{j\to\infty}\mu_{g_j}(E_j)=v.
\]
Moreover,
\[
 \Per_g(E_j)\leq e^{(n-1)\delta_j}\Per_{g_j}(E_j).
\]
The right-hand side is bounded by the limsup construction.  Hence lower
semicontinuity and $\delta_j\to0$ give
\[
 \cI_g(v)\leq\Per_g(E)
 \leq\liminf_{j\to\infty}\Per_g(E_j)
 \leq\liminf_{j\to\infty}\Per_{g_j}(E_j).
\]
Together with \eqref{eq:stability-limsup}, this proves the asserted convergence
\eqref{eq:profile-convergence}.  The same compactness argument applied to any
sequence of minimisers produces a limit $E$ for which equality holds
throughout; hence $E$ is minimising and
\eqref{eq:minimiser-perimeter-convergence} follows.
\end{proof}

For $q\in\Met(M)$ define the normalised profile
\[
 \widehat{\cI}_q(\theta):=\cI_q(\theta V_q),
 \qquad 0\leq\theta\leq1.
\]

\begin{corollary}[Quantitative stability of normalised profiles]
\label{cor:normalized-profile-holder}
Fix $g\in\Met(M)$ and $0<\varepsilon<1/2$.  There are
$\delta_*>0$ and $C<\infty$ such that, whenever
$\delta:=\dMet(g,\widetilde g)\leq\delta_*$,
\begin{equation}\label{eq:normalized-profile-holder}
 \sup_{\theta\in[\varepsilon,1-\varepsilon]}
 \bigl|\widehat{\cI}_{\widetilde g}(\theta)
       -\widehat{\cI}_g(\theta)\bigr|
 \leq C\delta^{(n-1)/n}.
\end{equation}
Thus convergence of rough metrics gives locally uniform convergence of their
normalised isoperimetric profiles, with an explicit H\"older exponent.
\end{corollary}

\begin{proof}
Fix a smooth background metric $h_0$ and $\delta_0>0$, and put
\[
 D:=\dMet(g,h_0)+\delta_0,
 \qquad
 \alpha:=\varepsilon e^{-2n\delta_0}V_g,
 \qquad
 \gamma:=\frac{n-1}{n}.
\]
Let $s_0,C_R$ be the constants in Lemma~\ref{lem:uniform-repair} for
$(h_0,D,\alpha)$.  Choose
$0<\delta_*\leq\min\{\delta_0,1\}$ so small that
\begin{equation}\label{eq:delta-star-choice}
 e^{n\delta_0}V_g(e^{n\delta_*}-1)\leq s_0.
\end{equation}

Fix $\theta\in[\varepsilon,1-\varepsilon]$, and let $E$ be a
$g$-isoperimetric region with $\mu_g(E)=\theta V_g$.  If
$J=d\mu_{\widetilde g}/d\mu_g$, then
$e^{-n\delta}\leq J\leq e^{n\delta}$ and
\[
 s:=\theta V_{\widetilde g}-\mu_{\widetilde g}(E)
   =\int_M(\theta-\one_E)(J-1)\dmu g.
\]
Consequently,
\begin{equation}\label{eq:metric-volume-error}
 |s|\leq V_g(e^{n\delta}-1).
\end{equation}
Both phases of $E$ have $\widetilde g$-volume at least
$e^{-n\delta}\varepsilon V_g\geq\alpha$.  The uniform repair lemma, with
$q=\widetilde g$, and the perimeter comparison therefore give
\begin{equation}\label{eq:profile-two-sided-a}
 \widehat{\cI}_{\widetilde g}(\theta)
 \leq e^{(n-1)\delta}\widehat{\cI}_g(\theta)
   +C_0(e^{n\delta}-1)^\gamma
\end{equation}
for a constant $C_0$ independent of $\theta$ and $\widetilde g$.

Conversely, begin with a $\widetilde g$-isoperimetric region
$\widetilde E$ of volume $\theta V_{\widetilde g}$ and repair its $g$-volume.
Its two phases have $g$-volume at least
\[
 e^{-n\delta}\varepsilon V_{\widetilde g}
 \geq\varepsilon e^{-2n\delta_0}V_g=\alpha,
\]
and the reverse volume error is bounded by
\[
 V_{\widetilde g}(e^{n\delta}-1)
 \leq e^{n\delta_0}V_g(e^{n\delta}-1)\leq s_0.
\]
The same argument yields
\begin{equation}\label{eq:profile-two-sided-b}
 \widehat{\cI}_g(\theta)
 \leq e^{(n-1)\delta}\widehat{\cI}_{\widetilde g}(\theta)
   +C_0(e^{n\delta}-1)^\gamma.
\end{equation}
Proposition~\ref{prop:profile-continuity} bounds
$\widehat{\cI}_g$ on $[\varepsilon,1-\varepsilon]$, and
\eqref{eq:profile-two-sided-a} then bounds
$\widehat{\cI}_{\widetilde g}$ there as well.  According as
$\widehat{\cI}_{\widetilde g}(\theta)\geq\widehat{\cI}_g(\theta)$ or the
reverse inequality holds, use \eqref{eq:profile-two-sided-a} or
\eqref{eq:profile-two-sided-b} to obtain
\[
 \bigl|\widehat{\cI}_{\widetilde g}(\theta)
       -\widehat{\cI}_g(\theta)\bigr|
 \leq C_1\left((e^{(n-1)\delta}-1)
                +(e^{n\delta}-1)^\gamma\right).
\]
Finally, $e^{a\delta}-1\leq C_a\delta$ for $0\leq\delta\leq1$ and
$\delta\leq\delta^\gamma$ on that interval prove
\eqref{eq:normalized-profile-holder}.
\end{proof}

\begin{remark}[A strong form of $\Gamma$-convergence]
The proof uses more than abstract $\Gamma$-convergence: the two-sided ratio
bounds in \eqref{eq:perimeter-comparison} tend to one uniformly in $E$.  This
does not assert a uniform additive bound on sets of unbounded perimeter.  The
only nontrivial point is repairing the small mismatch between the two volume
constraints.  This is why the entire stability theorem is short.
\end{remark}

\section{Isoperimetric constants, Cheeger sets, and the rough Laplacian}

Define the scale-invariant relative isoperimetric constant and the Cheeger
constant by
\begin{align}
 \cS(g)&:=\inf_{0<\mu_g(E)<V_g}
 \frac{\Per_g(E)}
 {\min\{\mu_g(E),\mu_g(M\setminus E)\}^{(n-1)/n}},
 \label{eq:S-constant}\\
 \hch(g)&:=\inf_{0<\mu_g(E)<V_g}
 \frac{\Per_g(E)}
 {\min\{\mu_g(E),\mu_g(M\setminus E)\}}.
 \label{eq:Cheeger-constant}
\end{align}
The lower bound \eqref{eq:rough-relative-iso} shows that both are positive.
The Cheeger infimum is attained: a minimising sequence cannot lose all of
one phase, because its quotient would diverge like
$\min\{\mu_g(E),\mu_g(E^c)\}^{-1/n}$; BV compactness and lower
semicontinuity then apply.  This existence mechanism is a special case of the
general relaxed-perimeter Cheeger theory; see, for example,
\cite{FranceschiPinamontiSaraccoStefani2024}.  What is useful here is its
placement in the Bandara--Hassan space, with the explicit stability constants
below.

\begin{theorem}[Log-Lipschitz geometric invariants]
\label{thm:geometric-log-Lipschitz}
For $g,\widetilde g\in\Met(M)$ and
$\delta=\dMet(g,\widetilde g)$,
\begin{align}
 e^{-2(n-1)\delta}\cS(g)
 &\leq\cS(\widetilde g)
 \leq e^{2(n-1)\delta}\cS(g),
 \label{eq:S-stability}\\
 e^{-(2n-1)\delta}\hch(g)
 &\leq\hch(\widetilde g)
 \leq e^{(2n-1)\delta}\hch(g).
 \label{eq:h-stability}
\end{align}
Equivalently, $\log\cS$ and $\log\hch$ are globally Lipschitz on the
space of rough metrics, with constants $2(n-1)$ and $2n-1$, respectively.
\end{theorem}

\begin{proof}
For every fixed set $E$, \eqref{eq:measure-comparison} compares the smaller
phase volume by factors $e^{\pm n\delta}$, while
\eqref{eq:perimeter-comparison} compares the numerator by factors
$e^{\pm(n-1)\delta}$.  Raise the volume comparison to the power
$(n-1)/n$ for \eqref{eq:S-stability}, and then take the infimum.  Without
that power the exponents add to $2n-1$, giving \eqref{eq:h-stability}.
\end{proof}

We now turn to the operator-theoretic side.  The closed form
\begin{equation}\label{eq:Dirichlet-form}
 \cE_g(u):=\int_M |du|_{g^*}^2\dmu g,
 \qquad u\in H^1(M,g),
\end{equation}
defines the nonnegative self-adjoint rough Laplacian $\Delta_g$.  On compact
$M$, comparison with a smooth metric gives compact resolvent, and hence
\[
 0=\lambda_0(g)<\lambda_1(g)\leq\lambda_2(g)\leq\cdots\nearrow\infty,
\]
with multiplicity.  This construction and much more of the analytic theory
are developed in
\cite{Bandara2016,BandaraBryan2020,BandaraNursultanovRowlett2021}.
We stress that $\cE_g$ uses the full measurable coefficient tensor.  It need
not coincide with the metric Cheeger energy of the pair consisting of Lashi's
intrinsic distance and $\mu_g$; that pair can fail to recover the tensor
\cite{BandaraHassan2025}.  For the same reason, $\Per_g$ is the relaxed rough
perimeter defined in \eqref{eq:rough-TV}, not a naively induced hypersurface
area or an unqualified metric-measure perimeter.

\begin{theorem}[Every eigenvalue is stable]
\label{thm:eigenvalue-stability}
For $g,\widetilde g\in\Met(M)$, set
$\delta=\dMet(g,\widetilde g)$.  Then, for every $k\geq1$,
\begin{equation}\label{eq:eigenvalue-comparison}
 e^{-2n\delta}\lambda_k(g)
 \leq\lambda_k(\widetilde g)
 \leq e^{2n\delta}\lambda_k(g).
\end{equation}
Thus $g\mapsto\log\lambda_k(g)$ is $2n$-Lipschitz on
$(\Met(M),\dMet)$, uniformly in $k$.
\end{theorem}

\begin{proof}
The form domains for $g$ and $\widetilde g$ agree as sets.  Lemma
\ref{lem:cofactor}, with $p=2$, gives
\[
 e^{-n\delta}\cE_g(u)\leq\cE_{\widetilde g}(u)
 \leq e^{n\delta}\cE_g(u),
\]
while \eqref{eq:measure-comparison} gives
\[
 e^{-n\delta}\|u\|_{L^2(g)}^2
 \leq\|u\|_{L^2(\widetilde g)}^2
 \leq e^{n\delta}\|u\|_{L^2(g)}^2.
\]
Therefore every Rayleigh quotient is compared by factors
$e^{\pm2n\delta}$.  Apply the min--max characterisation of
$\lambda_k$.
\end{proof}

The geometric and operator-theoretic quantities are related exactly as in the
smooth setting.

\begin{lemma}[A coarea inequality sufficient for Cheeger]
\label{lem:weak-coarea}
If $w$ is a nonnegative Lipschitz function on $M$, then
\begin{equation}\label{eq:weak-coarea}
 \int_0^\infty\Per_g(\{w>t\})\,dt
 \leq\int_M|dw|_{g^*}\dmu g.
\end{equation}
\end{lemma}

\begin{proof}
Choose a smooth nondecreasing $\eta:\R\to[0,1]$ which is zero on
$(-\infty,0]$, one on $[1,\infty)$, and satisfies
$\int_\R\eta'=1$.  For every $t$, as $\varepsilon\downarrow0$,
\[
 \eta\left(\frac{w-t}{\varepsilon}\right)
 \longrightarrow\one_{\{w>t\}}
 \quad\text{in }L^1.
\]
The displayed truncations are Lipschitz.  They may be approximated in
$W^{1,1}$ by smooth functions; comparison with a smooth background metric
and the Lipschitz estimate for $F_{g,h}$ show that their rough energies also
converge.
Consequently, with
$v_{\varepsilon,t}:=\eta((w-t)/\varepsilon)$,
\[
 \TV_g(v_{\varepsilon,t})
 \leq\int_M\frac1\varepsilon
 \eta'\left(\frac{w-t}{\varepsilon}\right)|dw|_{g^*}\dmu g.
\]
As $\varepsilon\downarrow0$, these functions converge in $L^1$ to
$\one_{\{w>t\}}$.  Lower semicontinuity gives
$\Per_g(\{w>t\})\leq\liminf_{\varepsilon\downarrow0}
\TV_g(v_{\varepsilon,t})$.  Fatou's lemma and Fubini's theorem therefore yield
\begin{align*}
 \int_0^\infty\Per_g(\{w>t\})\,dt
 &\leq\liminf_{\varepsilon\downarrow0}
 \int_\R\int_M
 \frac1\varepsilon\eta'\left(\frac{w-t}{\varepsilon}\right)
 |dw|_{g^*}\dmu g\,dt\\
 &=\int_M|dw|_{g^*}\dmu g.
\end{align*}
\end{proof}

\begin{theorem}[Cheeger's inequality for a rough metric]
\label{thm:rough-Cheeger}
For every rough metric on compact connected $M$,
\begin{equation}\label{eq:rough-Cheeger}
 \lambda_1(g)\geq\frac{\hch(g)^2}{4}.
\end{equation}
\end{theorem}

\begin{proof}
Let $u\in C^\infty(M)$ and choose a median $m$ of $u$, so that both
$\{u>m\}$ and $\{u<m\}$ have measure at most $V_g/2$.  Put
\[
 u_+=(u-m)_+,
 \qquad
 u_-=(m-u)_+.
\]
If $f$ is either $u_+$ or $u_-$, apply
Lemma~\ref{lem:weak-coarea} to $w=f^2$.  Every positive superlevel set of
$w$ occupies at most half the volume, so the definition of $\hch(g)$ and
the layer-cake formula give
\begin{align*}
 \hch(g)\int_M f^2\dmu g
 &\leq\int_0^\infty\Per_g(\{f^2>t\})\,dt\\
 &\leq2\int_M f|df|_{g^*}\dmu g\\
 &\leq2\left(\int_M f^2\dmu g\right)^{1/2}
       \left(\int_M|df|_{g^*}^2\dmu g\right)^{1/2}.
\end{align*}
Consequently,
\[
 \int_M|df|_{g^*}^2\dmu g
 \geq\frac{\hch(g)^2}{4}\int_Mf^2\dmu g.
\]
The Sobolev chain rule gives
$du_+=\one_{\{u>m\}}du$, $du_-=-\one_{\{u<m\}}du$, and $du=0$ almost
everywhere on $\{u=m\}$.  Adding the estimates for $u_+$ and $u_-$ therefore
gives
\[
 \cE_g(u)\geq\frac{\hch(g)^2}{4}
 \int_M|u-m|^2\dmu g
 \geq\frac{\hch(g)^2}{4}
 \inf_{c\in\R}\int_M|u-c|^2\dmu g.
\]
Smooth density and the variational identity
\[
 \lambda_1(g)=\inf_{u\not\equiv\mathrm{const}}
 \frac{\cE_g(u)}{\displaystyle
  \inf_{c\in\R}\int_M|u-c|^2\dmu g}
\]
prove \eqref{eq:rough-Cheeger}.
\end{proof}

\section{What this says about rough CMC}

For every rough metric and every prescribed
volume, there is a rough isoperimetric region---a variational CMC analogue.
These regions are stable
under perturbation of the rough metric in exactly the topology introduced by
Bandara and Hassan.  No Christoffel symbols and no derivatives of $g$ enter.

When $g$ is smooth, standard isoperimetric regularity
\cite{Maggi2012} says that the reduced boundary of $E_v$ is a smooth embedded
CMC hypersurface away from a singular set of Hausdorff dimension at most
$n-8$.  Thus the singular set is empty when
$2\leq n\leq7$.  In this regime our variational definition returns the usual
CMC equation
\[
 H_g=\lambda
\]
on the regular boundary, with $\lambda$ the volume-constraint multiplier.

For a genuinely rough metric we should resist writing this equation too
early.  The geometric question is instead a regularity ladder.

\begin{enumerate}
\item \emph{Measurable coefficients.}  Existence, compactness, coarea,
  Cheeger theory, and spectral comparison survive.  Classical hypersurface
  regularity need not survive, and the relaxed perimeter may encode a
  nontrivial effective surface tension.
\item \emph{VMO, Dini, or H\"older coefficients.}  Determine the weakest
  modulus under which excess decay and partial regularity hold for
  $\Per_g$-minimisers.  This is the natural interface with the regularity
  theory of elliptic parametric integrals.
\item \emph{Differentiable coefficients.}  Identify a threshold at which the
  first variation exists in a sufficiently strong form, recover a weak CMC
  equation, and bootstrap it to the regularity permitted by $g$.
\end{enumerate}

A minimising movement scheme could be used to construct a rough volume-preserving mean-curvature flow, perhaps in another paper.

\section{Some ideas}

The results above can form the core of a  first paper.
Here are some more things to do.

\begin{enumerate}
\item Write the manifold version of the Amar--Bellettini dual
  representation explicitly.  The relaxation definition
  \eqref{eq:rough-TV} is already invariant and sufficient for all the
  theorems above; the dual formula will clarify calibrations and is a
  prerequisite for formulating a weak mean-curvature theory.
\item Decide the first useful regularity threshold.  The most efficient target
  is probably $g\in C^{0,\alpha}$ for partial regularity, followed by
  $g\in C^{1,\alpha}$ for the classical CMC equation.  The merely measurable
  theorem should remain as the compactness base case.
\item Study uniqueness and quantitative stability near a nondegenerate smooth
  isoperimetric region.  Theorem~\ref{thm:profile-stability} gives compactness;
  an implicit-function or second-variation argument should upgrade this to
  convergence without subsequences and, under a spectral gap, a rate.
\item Then return to the Almgren--Taylor--Wang step.  Use the relaxed
  perimeter as the energy and Lashi's intrinsic distance for the mobility
  term.  Direct-method and almost-minimality arguments should then go through with the geometrically correct energy.
\end{enumerate}

My main point is that we already have an interesting theorem: rough isoperimetric regions exist and vary continuously with the
rough metric, while the profiles vary H\"older-continuously and the associated
geometric and spectral bottleneck
quantities vary quantitatively.  The remaining CMC problem is a precise and
natural regularity problem rather than an existence problem.

\begin{thebibliography}{99}

\bibitem{AmarBellettini1994}
M.~Amar and G.~Bellettini,
\newblock A notion of total variation depending on a metric with discontinuous
coefficients,
\newblock \emph{Ann. Inst. H. Poincar\'e Anal. Non Lin\'eaire} \textbf{11}
(1994), no.~1, 91--133.

\bibitem{AFP2000}
L.~Ambrosio, N.~Fusco and D.~Pallara,
\newblock \emph{Functions of bounded variation and free discontinuity
problems},
\newblock Oxford Mathematical Monographs, Oxford University Press, 2000.

\bibitem{Bandara2016}
L.~Bandara,
\newblock Rough metrics on manifolds and quadratic estimates,
\newblock \emph{Math. Z.} \textbf{283} (2016), no.~3--4, 1245--1281.

\bibitem{BandaraBryan2020}
L.~Bandara and P.~Bryan,
\newblock Heat kernels and regularity for rough metrics on smooth manifolds,
\newblock \emph{Math. Nachr.} \textbf{293} (2020), no.~12, 2255--2270.

\bibitem{BandaraHassan2025}
L.~Bandara and A.~Hassan,
\newblock The space of rough Riemannian metrics,
\newblock preprint, \href{https://arxiv.org/abs/2507.10274}{arXiv:2507.10274}
(2025).

\bibitem{BandaraNursultanovRowlett2021}
L.~Bandara, M.~Nursultanov and J.~Rowlett,
\newblock Eigenvalue asymptotics for weighted Laplace equations on rough
Riemannian manifolds with boundary,
\newblock \emph{Ann. Sc. Norm. Super. Pisa Cl. Sci. (5)} \textbf{22} (2021),
no.~4, 1843--1878.

\bibitem{BellettiniBouchitteFragala1999}
G.~Bellettini, G.~Bouchitt\'e and I.~Fragal\`a,
\newblock BV functions with respect to a measure and relaxation of metric
integral functionals,
\newblock \emph{J. Convex Anal.} \textbf{6} (1999), no.~2, 349--366.

\bibitem{Cheeger1970}
J.~Cheeger,
\newblock A lower bound for the smallest eigenvalue of the Laplacian,
\newblock in \emph{Problems in analysis}, Princeton University Press, 1970,
195--199.

\bibitem{FranceschiPinamontiSaraccoStefani2024}
V.~Franceschi, A.~Pinamonti, G.~Saracco and G.~Stefani,
\newblock The Cheeger problem in abstract measure spaces,
\newblock \emph{J. Lond. Math. Soc. (2)} \textbf{109} (2024), no.~1,
Paper No.~e12840.

\bibitem{Maggi2012}
F.~Maggi,
\newblock \emph{Sets of finite perimeter and geometric variational problems},
\newblock Cambridge Studies in Advanced Mathematics, vol.~135, Cambridge
University Press, 2012.

\bibitem{Zheng2020}
Y.~Zheng,
\newblock Uniform Lipschitz continuity of the isoperimetric profile of compact
surfaces under normalized Ricci flow,
\newblock preprint, \href{https://arxiv.org/abs/2001.00341}{arXiv:2001.00341}
(2020).

\end{thebibliography}

\end{document}
